\begin{resumo}
Esta dissertação investiga como técnicas de avaliação de sistemas de geração aumentada por recuperação (RAG) podem ser organizadas e aplicadas à análise da recuperação documental, da geração de respostas e da verificabilidade das fontes em um sistema voltado à consulta de discursos da 56\textordfeminine{} Legislatura do Senado Federal. A pesquisa é aplicada, exploratória e descritiva, adota abordagem multimétodos e dialoga com a \textit{Design Science Research}. O corpus reúne 15.729 registros de pronunciamentos; 15.726 contêm texto integral ou resumo utilizável. Para o percurso de viabilidade, foi elaborada uma bateria de 200 perguntas em 16 categorias e selecionado um piloto estratificado de 32 itens. Na execução reportada, todos os itens foram processados, e o juiz automatizado atribuiu média de 9,094 em escala de 0 a 10, segundo uma rubrica exploratória. O registro também apontou alto risco de mistura de autores em 28 itens. Essa média não constitui avaliação humana independente: \texttt{gpt-5-nano} foi usado tanto para gerar quanto para julgar as respostas. O experimento principal ainda não foi realizado; não há conjunto de referência validado, comparação controlada entre recuperadores ou avaliação humana independente.
\end{resumo}
